\documentclass[12pt]{article}
\usepackage{amsmath, amssymb, amsthm}
\usepackage{geometry}
\geometry{margin=1in}

\title{CSE400: Lecture 21 Scribe \\ Inequalities and Tail Bounds}
\author{}
\date{}

\begin{document}

\maketitle

\section{Introduction to Tails}

\subsection{Definition}
The \textbf{tail} of a random variable $X$ is:
\[
P(X > x)
\]

\subsection{Motivation}
We care about tail probabilities because they represent \textbf{rare extreme events}:
\begin{itemize}
    \item Jobs delayed more than 24 hours
    \item Packets dropped due to full buffer
\end{itemize}

\textbf{Key Idea:}
\begin{itemize}
    \item Mean $\rightarrow$ average behavior (easy)
    \item Tail $\rightarrow$ rare extreme events (hard)
\end{itemize}

\textbf{Why hard?}
\begin{itemize}
    \item Requires summing many small probabilities
\end{itemize}

---

\section{Tail Examples}

\subsection{Example 1: Load Balancing}
Suppose $n$ jobs are distributed among $n$ servers.

Let $X =$ number of jobs assigned to a server.

\[
X \sim \text{Binomial}(n,p)
\]

We want:
\[
P(X \ge k) = \sum_{i=k}^{n} \binom{n}{i} p^i (1-p)^{n-i}
\]

This represents a \textbf{tail event} (rare overload).

---

\subsection{Example 2: Poisson Arrivals}

Jobs arrive with rate $\lambda$.

\[
X \sim \text{Poisson}(\lambda)
\]

\[
P(X \ge k) = \sum_{i=k}^{\infty} e^{-\lambda} \frac{\lambda^i}{i!}
\]

This represents a \textbf{sudden spike in load}.

---

\section{Why Tail Bounds?}

Exact computation:
\[
P(X \ge k) = \sum \text{(long sums)}
\]

\textbf{Problem:} Computationally complex.

\textbf{Idea:}
\begin{itemize}
    \item Find an \textbf{upper bound}
    \item Easier to compute
    \item Still guarantees correctness
\end{itemize}

\[
P(X \ge k) \le \text{simple bound}
\]

\textbf{Key Idea:} Approximate safely instead of computing exactly.

---

\section{Running Example}

Flip a fair coin $n$ times.

\[
X \sim \text{Binomial}(n, \tfrac{1}{2})
\]

\[
E[X] = \frac{n}{2}
\]

Question:
\[
P\left(X \ge \frac{3n}{4}\right)
\]

---

\section{Markov's Inequality}

\subsection{Statement}

For non-negative $X$:
\[
P(X \ge a) \le \frac{E[X]}{a}
\]

\subsection{Key Idea}
\begin{itemize}
    \item Uses only mean
    \item Large values are rare
\end{itemize}

---

\subsection{Proof}

\[
E[X] = \sum_{x} x \cdot p_X(x)
\]

Consider only $x \ge a$:

\[
E[X] \ge \sum_{x \ge a} x p_X(x)
\]

Since $x \ge a$:

\[
E[X] \ge \sum_{x \ge a} a \cdot p_X(x)
= a \sum_{x \ge a} p_X(x)
= a P(X \ge a)
\]

Thus:
\[
P(X \ge a) \le \frac{E[X]}{a}
\]

---

\subsection{Running Example Application}

\[
E[X] = \frac{n}{2}, \quad a = \frac{3n}{4}
\]

\[
P\left(X \ge \frac{3n}{4}\right) \le \frac{n/2}{3n/4} = \frac{2}{3}
\]

\textbf{Observation:} Very loose bound (does not depend on $n$).

---

\subsection{Real-Life Example}

Average income = 10,000.

\[
P(X \ge 50,000) \le \frac{10,000}{50,000} = 0.2
\]

At most 20\% earn 50,000.

---

\subsection{Solved Example 1}

Given:
\[
E[X] = 70, \quad a = 90
\]

\[
P(X \ge 90) \le \frac{70}{90} = \frac{7}{9}
\]

---

\subsection{Solved Example 2}

\[
E[X] = 5000, \quad a = 20000
\]

\[
P(X \ge 20000) \le \frac{5000}{20000} = \frac{1}{4}
\]

---

\section{Chebyshev's Inequality}

\subsection{Key Idea}
\begin{itemize}
    \item Uses mean + variance
    \item Controls deviation from mean
\end{itemize}

\subsection{Statement}

\[
P(|X - \mu| \ge a) \le \frac{\text{Var}(X)}{a^2}
\]

---

\subsection{Proof Idea}

Apply Markov on $(X - \mu)^2$:

\[
P(|X - \mu| \ge a) = P((X-\mu)^2 \ge a^2)
\]

\[
\le \frac{E[(X-\mu)^2]}{a^2}
= \frac{\text{Var}(X)}{a^2}
\]

---

\subsection{Running Example}

\[
\mu = \frac{n}{2}, \quad \text{Var}(X) = \frac{n}{4}
\]

\[
P\left(X \ge \frac{3n}{4}\right)
= P\left(X - \frac{n}{2} \ge \frac{n}{4}\right)
\]

\[
\le \frac{n/4}{(n/4)^2} = \frac{1}{n}
\]

\textbf{Better than Markov (depends on $n$).}

---

\subsection{Real-Life Example}

\[
\mu = 70, \quad \sigma = 10
\]

\[
P(|X - 70| \ge 30) \le \frac{100}{900} = \frac{1}{9}
\]

---

\subsection{Solved Example 1}

\[
\mu = 70, \quad \sigma^2 = 25, \quad a = 20
\]

\[
P(X \ge 90) \le P(|X-70| \ge 20)
\]

\[
\le \frac{25}{400} = \frac{1}{16}
\]

---

\subsection{Solved Example 2}

\[
\mu = 50, \quad \sigma^2 = 16, \quad a = 20
\]

\[
P(|X-50| \ge 20) \le \frac{16}{400} = \frac{1}{25}
\]

---

\section{Chernoff Bound}

\subsection{Key Idea}

Progression:
\begin{itemize}
    \item Markov: use $X$
    \item Chebyshev: use $(X-\mu)^2$
    \item Chernoff: use $e^{tX}$
\end{itemize}

\textbf{Insight:}
\begin{itemize}
    \item Exponential grows very fast
    \item Amplifies large values
    \item Leads to tighter bounds
\end{itemize}

---

\subsection{Core Trick}

\[
P(X \ge a) = P(e^{tX} \ge e^{ta})
\]

Apply Markov:

\[
P(e^{tX} \ge e^{ta}) \le \frac{E[e^{tX}]}{e^{ta}}
\]

---

\subsection{Final Form}

\[
P(X \ge a) \le \min_{t > 0} \frac{E[e^{tX}]}{e^{ta}}
\]

\textbf{Key Insight:}
\begin{itemize}
    \item Holds for all $t > 0$
    \item Choose best $t$ for tightest bound
\end{itemize}

---

\subsection{Lower Tail}

\[
P(X \le a) \le \min_{t < 0} \frac{E[e^{tX}]}{e^{ta}}
\]

---

\section{Chernoff for Poisson}

\subsection{Goal}

Bound tail for:
\[
X \sim \text{Poisson}(\lambda)
\]

\subsection{Step 1: Compute MGF}

\[
E[e^{tX}] = \sum_{i=0}^{\infty} e^{ti} \cdot e^{-\lambda} \frac{\lambda^i}{i!}
\]

\[
= e^{-\lambda} \sum_{i=0}^{\infty} \frac{(\lambda e^t)^i}{i!}
\]

\[
= e^{-\lambda} e^{\lambda e^t}
= e^{\lambda (e^t - 1)}
\]

---

\textbf{This is used inside Chernoff bound to get final tail bounds.}

---

\end{document}